\begin{table}[h]
\centering
\caption{NVE threshold sensitivity analysis for CIVO detection. Task~B2 (main text) uses a 1.5$\times$ threshold yielding one false positive (008e\_v3\_s2, identified as non-CIVO Imagination-Reality Divergence). Task~B1 uses a stricter 1.7$\times$ threshold achieving perfect separation but with only 2 positive samples. All CIs are two-sided Clopper-Pearson exact binomial ($\alpha/2 = 0.025$).}
\label{tab:nve_sensitivity}
\small
\setlength{\tabcolsep}{4pt}
\begin{tabular}{@{}lcccccccl@{}}
\toprule
Task & Threshold & $n$ & Pos/Neg & AUROC & Accuracy & CP 95\% CI & Margin & Note \\
\midrule
B2 (main) & $\geq 1.5\times$ & 8 & 3\,/\,5 & 0.933 & 87.5\% & [0.474, 0.997] & $-2.16$ & 1 FP (008e\_v3\_s2: IRD) \\
B1 (sens.) & $\geq 1.7\times$ & 8 & 2\,/\,6 & 1.000 & 100\% & [0.631, 1.000] & $+0.35$ & $n_{\text{pos}}=2$, low power \\
\bottomrule
\end{tabular}

\vspace{4pt}
\noindent\small\textit{The stricter 1.7$\times$ threshold (B1) achieves perfect separation but with only $n_{\text{pos}}=2$, limiting statistical power. The 1.5$\times$ threshold (B2) increases recall at the cost of one false positive driven by non-CIVO Imagination-Reality Divergence.}
\end{table}
